% Cote 119: offprints, reports, courses, CV and letters (typed and printed part): archive cote 119, site page 225 (PDF page 226).
% typescript. Draft compared: none.
% AI-transcribed by Claude Opus 5.5 (claude-opus-5-5, adaptive thinking, effort high), Message Batches API, batch r119_full (msgbatch_01DwZNm5Y7YsKBWyAN4sFs8w), 2026-11-13; not yet checked by eye.
\scanpage{119/225}
\mspage{. 12 .}

\uline{Références} : J.P. Serre, A. Weil, B. Dwork, J. Tate,
M. Artin, P. Deligne ...

10. \uline{Géométrie Algébrique} ([12] à [19],[15 bis] à [20 bis]).

C'est dans cette direction que mon influence
a été la plus directe et la plus profonde, puisque
c'est dans cette optique que se placent pour l'essentiel
mes travaux depuis 1959. Voici les thèmes principaux
sous lesquels on peut placer mes contributions :

a) \uline{Travail de fondement} : Il s'agissait de dégager
un cadre suffisamment vaste pour servir de fondement
commun à la géométrie algébrique habituelle (y compris
celle développée par des auteurs comme A. Weil,
O. Zariski, C. Chevalley, J.P. Serre sur des corps de
base quelconques) et à l'arithmétique. C'est fait pour
l'essentiel dans [15, Chap. I,II et des parties des
Ch. III et IV], avec l'introduction et l'étude de la
\uline{notion de schéma}. Des généralisations ont été développées
par la suite, dans le même esprit, avec les schémas
formels [15, Chap. I, par. 10], la théorie des
``algebraic spaces'' de M. Artin (cf. Knutson [$*$]), les
``algebraic stacks'' ou ``multiplicités algébriques'' de
P. Deligne et D. Mumford ($*$), des ``schémas relatifs''
de la thèse de M. Hakim [$*$] (en attendant les
``multiplicités formelles'' et les ``multiplicités algébriques
relatives'' sur des topos annelés généraux, etc). Ces
généralisations montrent la part conceptuelle importante
qui revient, dans le langage des schémas, à la notion
générale de la localisation, c'est à dire à celle de
\uline{topos} (dont il a été question au par. 3). Les fondements
développés dans [15] et [16] sont aujourd'hui le ``pain
quotidien'' de la grande majorité des géomètres algébristes,

\margin{225}
